\section{The geometric reduction and the finite-time obstruction}
\label{geo:section}
\label{sec:geometry}

We use scalar curvature in the convention
\(S(\omega)=\operatorname{tr}_{\omega}\operatorname{Ric}(\omega)\).
Let \(\omega_*\) be the Fubini--Study form on \(\CP^n\) with affine
potential \(\log(1+|Z|^2)\).  The geometric identities below hold for every
\(n\geq1\); the construction will use \(n=10\).
The radial and momentum descriptions follow the framework of
Calabi~\cite{Calabi1979} and Hwang--Singer~\cite{HwangSinger2002};
the Legendre-coordinate viewpoint is also central to the toric formulas
of Guillemin and Abreu~\cite{Guillemin1994,Abreu1998}. We give the
identities directly to fix the scalar convention and to verify
smoothness at the point orbit as well as at the divisor.
Set
\begin{equation}\label{geo:reference}
 q=|Z|^2,\qquad \ell=\log q,\qquad
 v_*(x)=x(1-x),\qquad w(x)=x^{n-1}.
\end{equation}
For a radial potential \(\Phi(\ell)\), its momentum coordinate and
momentum profile are
\[
 x=\Phi'(\ell),\qquad v(x)=\Phi''(\ell).
\]
The Legendre transform \(u(x)=x\ell-\Phi(\ell)\) satisfies
\(u'=\ell\) and \(u''=1/v\).  The reference transform is
\[
 u_*(x)=x\log x+(1-x)\log(1-x),\qquad u_*''=v_*^{-1}.
\]
Primes on functions of \(x\) will denote differentiation in \(x\).

\begin{proposition}[Exact fixed-class reduction]\label{geo:reduction}
Let \(I\) be an open time interval, and let
\(H\in C^\infty(I\times[0,1])\) satisfy
\begin{equation}\label{geo:positivity}
 b(t,x):=1+v_*(x)H(t,x)>0.
\end{equation}
Define
\begin{equation}\label{geo:momentum}
 v=v_*+v_*^2H=v_*b,\qquad
 K(H)=w^{-1}(wv_*^2H)'',\qquad
 \A H=\big[w^{-1}(wv_*^2H)''\big]''.
\end{equation}
At each time these data, together with an affine choice in the Legendre
potential, determine a smooth \(U(n)\)-invariant K\"ahler form in
\([\omega_*]\).  Its scalar curvature and average scalar curvature are
\begin{equation}\label{geo:scalar}
 S(t,x)=n(n+1)-K(H)(t,x),\qquad \bar S=n(n+1).
\end{equation}
If, in addition,
\begin{equation}\label{geo:equation}
 H_t+(1+v_*H)^2\A H=0,
\end{equation}
then for any \(t_0\in I\) and any initial affine choice there is a smooth
family of these forms satisfying the Calabi flow exactly.  More
precisely, its global potential relative to \(\omega(t_0)\) satisfies
\begin{equation}\label{geo:potential-flow}
 \phi(t_0)=0,\qquad
 \omega(t)=\omega(t_0)+i\partial\bar\partial\phi(t),\qquad
 \phi_t=S(\omega(t))-\bar S.
\end{equation}
\end{proposition}

\begin{proof}
We first verify smoothness and positivity at both endpoints.  Suppress
time for this part of the argument.  Since
\begin{equation}\label{geo:smooth-legendre}
 \frac1v-\frac1{v_*}=-\frac{H}{1+v_*H},
\end{equation}
we can choose \(\psi\in C^\infty[0,1]\) with second derivative equal to
the right-hand side, and put \(u=u_*+\psi\).  The two integration
constants are precisely its free affine part.  Write \(a=\psi'\).
The function \(u'\) is strictly increasing from \(-\infty\) to
\(+\infty\), so it has a smooth inverse on the open interval.  In terms
of \(q=\exp\ell\), the inverse relation is
\begin{equation}\label{geo:center-inverse}
 q=\frac{x}{1-x}\exp a(x).
\end{equation}
At \(x=0\) its right-hand side vanishes and has derivative
\(\exp a(0)>0\).  The inverse therefore extends smoothly to \(q=0\),
with
\[
 x(q)=\exp(-a(0))q+O(q^2).
\]
Here and below a smooth remainder is meant: for example the last error
is \(q^2\) times a smooth function of \(q\).  Since
\(\dd\Phi/\dd q=x(q)/q\), integration shows that the Legendre dual
extends smoothly as a function of \(q\), with expansion
\begin{equation}\label{geo:center-expansion}
 \Phi(\log q)=C_0+\exp(-a(0))q+O(q^2).
\end{equation}
Division by \(q\) here preserves smoothness because \(x(0)=0\), as is
also clear from \(x(q)/q=\int_0^1x'(sq)\,\dd s\).
Its complex Hessian at \(Z=0\) is
\(\exp(-a(0))\) times the Euclidean Hermitian form, and is positive.

At the other endpoint set \(\eta=q^{-1}\) and \(b_1=1-x\).  The same
inverse relation becomes
\[
 \eta=\frac{b_1}{1-b_1}\exp\big(-a(1-b_1)\big),\qquad
 b_1(\eta)=\exp a(1)\eta+O(\eta^2).
\]
If \(\Psi(\eta)=\Phi(\log q)-\log q\), then
\(\Psi'(\eta)=b_1(\eta)/\eta\).  Consequently
\begin{equation}\label{geo:infinity-expansion}
 \Psi(\eta)=C_1+\exp a(1)\eta+O(\eta^2)
\end{equation}
is smooth at \(\eta=0\).  This expansion gives smoothness on the
projective manifold, not merely along radial rays.  Indeed, on the
projective chart meeting the divisor at infinity given by
\[
 W_0=1/Z_j,\qquad W_i=Z_i/Z_j\quad(i\ne j),\qquad
 B=1+\sum_{i\ne j}|W_i|^2,
\]
we have \(q=B/|W_0|^2\).  After removing the pluriharmonic term
\(-\log|W_0|^2\) from \(\Phi\), a local potential is
\begin{equation}\label{geo:divisor-potential}
 \log B+\Psi\!\left(\frac{|W_0|^2}{B}\right).
\end{equation}
On \(W_0=0\) its tangential complex Hessian is the Fubini--Study metric
of the divisor, its normal coefficient is \(\exp a(1)/B>0\), and its
mixed coefficients vanish.  This proves positivity there; for \(n=1\)
only the normal coefficient is present.  On \(0<x<1\), direct
differentiation gives the tangential and radial eigenvalues
\(x/q\) and \(v/q\), respectively, so positivity holds everywhere.

The difference
\(\Phi(\log|Z|^2)-\log(1+|Z|^2)\) is a globally smooth function.
At the origin this follows from \eqref{geo:center-expansion}; in the
chart at infinity it equals
\(\Psi(\eta)-\log(1+\eta)\), which is smooth by
\eqref{geo:infinity-expansion}.  Thus its complex Hessian is a global
exact form and the K\"ahler class is \([\omega_*]\).  All inversions and
integrations above depend smoothly on time on compact subintervals of
\(I\), by the positive endpoint derivatives and the inverse function
theorem with parameters.

We next compute scalar curvature.  The eigenvalues already found give
\[
 \det(g_{i\bar j})=q^{-n}x^{n-1}v=q^{-n}wv.
\]
For any radial function \(f\), the complex Laplacian is
\begin{equation}\label{geo:complex-laplacian}
 \Delta_\omega f
 =\frac{f_{\ell\ell}}{v}+\frac{n-1}{x}f_\ell
 =w^{-1}\partial_x(wf_\ell),
\end{equation}
where \(\partial_\ell=v\partial_x\) in the last expression.
For \(f=-\log\det g\), this yields
\[
 f_\ell=n-\frac{(wv)'}w,\qquad
 S=\frac{n(n-1)}x-w^{-1}(wv)''.
\]
Substituting \(v=v_*+v_*^2H\) gives
\eqref{geo:scalar}.  The expression \(K(H)\) is smooth at zero:
\(wv_*^2H=x^{n+1}(1-x)^2H\), whose second derivative is divisible by
\(x^{n-1}\) as a smooth function.  It is also smooth at one.

For completeness, the average can be read directly from this formula.
The radial volume density is a constant multiple of \(w(x)\,\dd x\):
the Euclidean polar factor \(q^{n-1}\dd q\), multiplied by the
determinant, becomes \(wv\,\dd\ell=w\,\dd x\).  Moreover,
\[
 \int_0^1 K(H)w\,\dd x
 =\big[(wv_*^2H)'\big]_{x=0}^{x=1}=0,
\]
because \(wv_*^2H\) vanishes to order at least two at both endpoints.
This proves the asserted average, including when \(n=1\).

Finally, fix \(t_0\in I\), choose \(u(t_0)\) as above, and define at
fixed momentum coordinate
\begin{equation}\label{geo:legendre-flow}
 u(t,x)=u(t_0,x)+\int_{t_0}^{t}K(H)(s,x)\,\dd s.
\end{equation}
Equation \eqref{geo:equation} implies
\[
 \partial_t(1/v)=-\frac{v_*^2H_t}{v^2}
 =\A H=\partial_x^2K(H)=\partial_tu''.
\]
Their initial values agree, so \(u''=1/v\) for all time.  In
particular the preceding metric construction applies throughout \(I\).
Differentiating the Legendre dual
\(\Phi(t,\ell)=x(t,\ell)\ell-u(t,x(t,\ell))\) at fixed \(\ell\)
cancels the terms containing \(x_t\), and gives
\begin{equation}\label{geo:legendre-sign}
 \Phi_t=-u_t=-K(H)=S-\bar S.
\end{equation}
The difference between \(\Phi(t)\) and \(\Phi(t_0)\) is globally smooth
by the endpoint argument and therefore supplies \(\phi\) in
\eqref{geo:potential-flow}.

This also specifies the affine part of \(u\) during evolution.
Replacing a fixed \(u\) by \(u+\alpha x+\beta\) changes its dual to
\(\Phi(\ell-\alpha)-\beta\), corresponding to the projective
dilation \(Z\mapsto\exp(-\alpha/2)Z\) and a potential constant.
An arbitrary time-dependent choice of these terms would affect
\eqref{geo:legendre-sign}; the prescription
\eqref{geo:legendre-flow} fixes them and establishes that equation
without an additional affine term.
\end{proof}

\subsection{A regular auxiliary sphere and its flat limit}

The endpoints of the momentum interval can be handled by a smooth
equation on a closed auxiliary manifold.  Set
\begin{equation}\label{geo:sphere}
 m=2n+4,\qquad
 \Sigma=S^{m+5}\subset\R^m\times\R^6,\qquad x=|z|^2.
\end{equation}
We identify functions of \(x\) with
\(O(m)\times O(6)\)-invariant functions on \(\Sigma\).

\begin{lemma}[The compact and flat operators]\label{geo:operators}
This identification is a bijection between \(C^\infty[0,1]\) and the
smooth invariant functions on \(\Sigma\).  The pushforward of spherical
volume under \(x\) is a positive constant multiple of
\(wv_*^2\,\dd x\).  With the nonpositive Laplacian convention, put
\(L=\frac14\Delta_\Sigma\).  On invariant functions,
\begin{equation}\label{geo:sphere-operator}
 L=x(1-x)\partial_x^2+\big(n+2-(n+5)x\big)\partial_x,
 \qquad
 \A=(-L+n+3)(-L+2(n+2)).
\end{equation}
The right-hand expression defines a smooth elliptic operator of order
four on all of \(\Sigma\).

In the core coordinates
\((z,\theta)\mapsto(z,\sqrt{1-|z|^2}\,\theta)\), with
\(\theta\in S^5\), its Laplacian on functions independent of \(\theta\)
is
\begin{equation}\label{geo:local-laplacian}
 \Delta_\Sigma
 =\Delta_z-\sum_{i,j=1}^m z_i z_j\partial_{z_i z_j}
             -(m+5)\sum_{i=1}^m z_i\partial_{z_i}.
\end{equation}
The flat momentum operator obtained by replacing \(v_*\) by \(x\) in
\eqref{geo:momentum} satisfies
\begin{equation}\label{geo:flat-operator}
 \wh\A H
 :=\big[x^{1-n}(x^{n+1}H)''\big]''
 =\big(x\partial_x^2+(n+2)\partial_x\big)^2H
 =\frac1{16}\Delta_{\R^m}^2H,\qquad x=|z|^2.
\end{equation}
\end{lemma}

\begin{proof}
The only regularity issue in the identification occurs at the two
singular orbits.  Here is an elementary smoothness argument.  If
\(f(r)\) is smooth and even near zero, then the operator
\[
 \mathcal T f(r)=\frac{f'(r)}r
              =\int_0^1f''(sr)\,\dd s
\]
defines another smooth even function.  Thus, for \(g(x)=f(\sqrt x)\),
each derivative on \(x>0\) has the form
\(g^{(j)}(x)=2^{-j}(\mathcal T^jf)(\sqrt x)\), and extends
continuously to \(x=0\).  Induction gives smoothness up to zero of
\(g\) and all its derivatives.  The same argument holds with smooth
parameters.  Restricting a smooth invariant function on \(\Sigma\) to
\((re_1,\sqrt{1-r^2}\,\theta_0)\) gives such an even \(f\), and
invariance makes its value depend only on \(x=r^2\).  Interchanging the
two factors proves smoothness in \(1-x\) at the other endpoint.
Conversely a smooth function of \(x\) is smooth in these coordinates
at both orbits, as well as on their complement.

Write \(|z|=\sin\vartheta\), with \(0\leq\vartheta\leq\pi/2\).
The spherical metric and volume density are
\[
 \dd\vartheta^2+\sin^2\vartheta\,g_{S^{m-1}}
                 +\cos^2\vartheta\,g_{S^5},\qquad
 \sin^{m-1}\vartheta\cos^5\vartheta\,\dd\vartheta\,
       \dd\omega_{m-1}\,\dd\omega_5.
\]
Changing to \(x=\sin^2\vartheta\) gives density proportional to
\[
 M(x)\,\dd x:=x^{m/2-1}(1-x)^2\,\dd x
              =wv_*^2\,\dd x.
\]
The radial Laplacian computed from this metric gives the displayed
formula for \(L\), equivalently
\(L=M^{-1}\partial_x(v_*M\partial_x)\).

We verify the factorization by a polynomial argument that also checks
the endpoint terms.  Both \(L\) and \(\A\) preserve polynomial degree
and are symmetric for the inner product with density \(M\).  For
\(\A\), twice integrating by parts gives
\begin{equation}\label{geo:symmetry}
 \int_0^1 a(\A b)M\,\dd x
 =\int_0^1\frac{(Ma)''(Mb)''}{w}\,\dd x.
\end{equation}
There are no boundary terms: \(Ma\) and its first derivative vanish
at each endpoint, and \(w^{-1}(Mb)''\) is smooth there.  For \(L\)
symmetry follows from its divergence expression, since \(v_*M\)
vanishes at both ends.  Orthogonalizing the monomials therefore gives
common polynomial eigenfunctions.  On the polynomial of degree \(j\),
the leading coefficients give respectively the eigenvalues
\[
 -j(j+n+4),\qquad
 (j+1)(j+2)(j+n+2)(j+n+3).
\]
Indeed a symmetric degree-preserving operator maps the degree-\(j\)
orthogonal polynomial to a multiple of itself, because its inner
product with every lower-degree polynomial vanishes.  Finally,
\[
 \big(j(j+n+4)+n+3\big)
 \big(j(j+n+4)+2(n+2)\big)
 =(j+1)(j+n+3)(j+2)(j+n+2).
\]
This proves equality on every polynomial, hence equality of the
differential expressions: at an interior point polynomials realize
arbitrary jets through order four.  The sphere expression has leading
term \(\Delta_\Sigma^2/16\), proving the claimed ellipticity.

For \eqref{geo:local-laplacian}, the inverse metric in the \(z\)
coordinates is \((\delta_{ij}-z_i z_j)\) and the volume density is
\((1-|z|^2)^2\,\dd z\,\dd\omega_5\).  Applying the divergence formula
therefore gives
\[
 (1-|z|^2)^{-2}\partial_{z_i}
 \big((1-|z|^2)^2(\delta_{ij}-z_i z_j)\partial_{z_j}\big)
 =(\delta_{ij}-z_i z_j)\partial_{z_i z_j}
   -(m+5)z_i\partial_{z_i},
\]
with repeated indices summed.  This computation only requires
independence of \(\theta\).  Finally both one-dimensional expressions
for \(\wh\A\) in \eqref{geo:flat-operator} expand to
\[
 x^2H^{(4)}+2(n+3)xH^{(3)}+(n+2)(n+3)H'',
\]
and the radial Euclidean identity
\(\frac14\Delta_{\R^m}=x\partial_x^2+\frac m2\partial_x\)
proves the last equality.
\end{proof}

\subsection{The self-similar core}
\label{geo:core-section}

From now on set \(n=10\), so \(m=24\) and \(\Sigma=S^{29}\).  In the
flat model the equation for the coefficient is
\begin{equation}\label{geo:flat-flow}
 F_t+(1+xF)^2\wh\A F=0.
\end{equation}
Let \(D=x\partial_x\).  For the function \(g=1+xF\),
\eqref{geo:flat-operator} gives
\begin{equation}\label{geo:euler-identity}
 x\wh\A\big(x^{-1}(g-1)\big)
 =x^{-2}P(D)(g-1),\qquad
 P(D)=(D-2)(D-1)(D+9)(D+10).
\end{equation}
Here \(P\) denotes a polynomial, as distinct from the evolution
operator denoted by \(P\) later.  One can check
\eqref{geo:euler-identity} on \(g-1=x^j\): both sides equal
\((j-2)(j-1)(j+9)(j+10)x^{j-2}\); equality of the differential
expressions follows as in Lemma~\ref{geo:operators}.
Consequently \eqref{geo:flat-flow} is equivalent, for \(x>0\), to
\begin{equation}\label{geo:g-flow}
 g_t=-g^2x^{-2}P(D)(g-1).
\end{equation}

Define the cubic polynomial operator
\(Q(D)=(D-1)(D+9)(D+10)\).  Suppose that \(G\) is smooth and positive
on \([0,\infty)\), with \(G(0)=1\), and for some \(c>0\) satisfies
\begin{equation}\label{geo:profile-equation}
 Q(y\partial_y)(G-1)=\frac{y^2}{2}\left(\frac1G-\frac1c\right).
\end{equation}
Then set
\begin{equation}\label{geo:selfsimilar}
 \rho=\sqrt{-t},\qquad y=x/\rho,\qquad
 h(y)=\frac{G(y)-1}{y},\qquad
 F(t,x)=\rho^{-1}h(x/\rho),\qquad t<0.
\end{equation}
The value \(h(0)=G'(0)\) makes \(h\) smooth at zero, since
\(h(y)=\int_0^1G'(sy)\,\dd s\).  Thus \(F\) is smooth for \(t<0\)
including at the flat origin.  Applying \(D-2\), now with
\(D=y\partial_y\), to \eqref{geo:profile-equation} yields
\[
 P(D)(G-1)=-\frac{y^2DG}{2G^2}.
\]
On the other hand \(\rho_t=-1/(2\rho)\), so
\(\partial_tG(x/\rho)=DG/(2\rho^2)\).  These identities prove
\eqref{geo:g-flow}, and hence \eqref{geo:flat-flow}, first for \(x>0\)
and then at zero by smoothness.  The existence of the required
nonconstant positive \(G\) is the subject of
Theorem~\ref{prof:profile}.

The same expression \(F(t,x)\) is smooth on the compact sphere for
every \(t<0\).  Its compact positivity multiplier obeys the useful
identity
\begin{equation}\label{geo:approximate-positivity}
 1+v_*F=x+(1-x)G(x/\rho).
\end{equation}
In particular, if \(G\) is bounded above and below by positive
constants, this multiplier has the same property uniformly in
\(t<0\) and \(0\leq x\leq1\).
This is a bound in the momentum coordinate \(x\), whose relation
to the fixed affine coordinate varies under the Legendre transform.
It does not imply uniform equivalence to the Fubini--Study metric.

\subsection{Curvature at the concentrating point}

The final geometric implication uses only one value of the correction.
It is convenient to state it before constructing that correction.

\begin{lemma}[Finite-time obstruction]\label{geo:obstruction}
Let \(G\in C^\infty[0,\infty)\) satisfy \(G(0)=1\) and
\(G'(0)\ne0\), and define \(F\) by \eqref{geo:selfsimilar}.  Suppose
that for some \(T>0\) and \(\delta>0\) there is a smooth invariant
solution \(H=F+V\) of \eqref{geo:equation} on \((-T,0)\times\Sigma\)
such that \(1+v_*H>0\) and
\begin{equation}\label{geo:correction-bound}
 V(t,0)=O(\rho^{-1+\delta})\qquad(t\uparrow0).
\end{equation}
Then a smooth \(U(10)\)-invariant metric \(\omega_0\in[\omega_*]\)
has maximal smooth Calabi-flow existence interval \([0,T_*)\), where
\(T_*=T/2<\infty\).  At the fixed affine origin \(p\), the scalar
curvature of this flow satisfies
\begin{equation}\label{geo:curvature-rate}
 S(\tau,p)=a(T_*-\tau)^{-1/2}
       +O\big(1+(T_*-\tau)^{(-1+\delta)/2}\big),\qquad
 a=-110G'(0)\ne0.
\end{equation}
In particular \(S(\tau,p)\sim a(T_*-\tau)^{-1/2}\), and \(a>0\)
when \(G'(0)<0\).
\end{lemma}

\begin{proof}
By Lemma~\ref{geo:operators}, \(H\) is smooth in \(x\) at both
endpoints for each time.  Proposition~\ref{geo:reduction}, with
\(t_0=-T/2\), gives a smooth family of K\"ahler forms, and its slice
\(\omega_0=\omega(-T/2)\) is a smooth positive form in the
Fubini--Study class.  Put \(\tau=t+T/2\).  The corresponding global
potential satisfies exactly the normalization and evolution in
\eqref{geo:potential-flow} on \(0\leq\tau<T_*=T/2\).

At \(x=0\), the coefficient of \(x^{n+1}\) in
\(wv_*^2H\) is \(H(t,0)\).  Thus \eqref{geo:scalar} gives the exact
identity
\begin{equation}\label{geo:center-scalar}
 S(t,p)=n(n+1)\big(1-H(t,0)\big).
\end{equation}
This evaluates scalar curvature at the same point \(p\) at every
time: \eqref{geo:center-inverse} sends \(q=0\) to \(x=0\) for every
affine choice.  Since
\[
 H(t,0)=\rho^{-1}G'(0)+O(\rho^{-1+\delta}),
\]
Equation~\eqref{geo:center-scalar} proves
\eqref{geo:curvature-rate} and the asserted divergence.

We recall why uniqueness applies to comparison with an unrestricted
smooth flow; no symmetry hypothesis is needed for this step.  Smooth
positive initial potentials have a unique local Calabi flow by
\cite[Theorem~3.2]{ChenHe2008}.  The uniqueness mechanism can also be
seen directly.  For two smooth solutions with the same initial
potential, let \(\chi\) be their difference and interpolate their
potentials by \(\varphi_s\), \(0\leq s\leq1\).  Their interpolated
K\"ahler forms remain positive, and variation of scalar curvature
gives
\[
 \chi_t=-\int_0^1
 \left(\Delta_{\varphi_s}^2\chi
       +\operatorname{Ric}_{\varphi_s}^{i\bar j}
                            \partial_i\partial_{\bar j}\chi\right)
 \,\dd s.
\]
On a compact time interval where both solutions are smooth, the
coefficients are smooth and uniformly elliptic.  Relative to a fixed
smooth volume form, integration twice by parts and interpolation give
\[
 \frac{\dd}{\dd t}\|\chi\|_{L^2}^2
 \leq-c\|\chi\|_{H^2}^2+C\|\chi\|_{L^2}^2
 \leq C\|\chi\|_{L^2}^2.
\]
For this estimate, the principal quadratic form is the integral of
\(|\Delta_{\varphi_s}\chi|^2\); differentiating the smooth volume
density introduces terms bounded by
\(C\|\chi\|_{H^1}\|\chi\|_{H^2}\), and the Ricci term has the same
or lower order.  The elliptic estimate for each uniformly elliptic
Laplacian and
\(\|\chi\|_{H^1}\leq\eps\|\chi\|_{H^2}+C_\eps\|\chi\|_{L^2}\)
absorb those terms.  Gronwall's inequality proves uniqueness.

It follows that any smooth flow from \(\omega_0\) agrees with the one
constructed above on their common interval.  Hence its maximal
interval contains \([0,T_*)\).  If it extended beyond \(T_*\), smooth
positivity on the compact manifold across that time would bound its
scalar curvature, contradicting \eqref{geo:curvature-rate}.  Its
maximal smooth existence time is therefore exactly \(T_*\).
\end{proof}

The Fubini--Study class contains the constant-scalar-curvature form
\(\omega_*\), so the
hypotheses of Lemma~\ref{geo:obstruction} suffice to refute the
smooth global-existence assertion.  Its requirement of smooth
existence at every finite time has already failed.  In addition,
\(\|S\|_{L^\infty}\) is unchanged by biholomorphic pullback, so the
curvature divergence itself is unaffected by allowing
automorphisms.
